\section{The Weil divisor class associated to an invertible module}
\label{section-c1}

\noindent
In this section we go through exactly the same progression as in
Section \ref{section-Weil-divisors} to define a canonical map
$\Pic(X) \to \text{Cl}(X)$
on a locally Noetherian integral scheme.

\medskip\noindent
Let $X$ be a scheme. Let $\mathcal{L}$ be an invertible $\mathcal{O}_X$-module.
Let $\xi \in X$ be a point. If $s_\xi, s'_\xi \in \mathcal{L}_\xi$ generate
$\mathcal{L}_\xi$ as $\mathcal{O}_{X, \xi}$-module, then there exists a unit
$u \in \mathcal{O}_{X, \xi}^*$ such that $s_\xi = u s'_\xi$.
The stalk of the sheaf of meromorphic sections
$\mathcal{K}_X(\mathcal{L})$ of $\mathcal{L}$
at $x$ is equal to
$\mathcal{K}_{X, x} \otimes_{\mathcal{O}_{X, x}} \mathcal{L}_x$.
Thus the image of any meromorphic section $s$
of $\mathcal{L}$ in the stalk at $x$ can be written as $s = fs_\xi$ with
$f \in \mathcal{K}_{X, x}$. Below we will abbreviate this by
saying $f = s/s_\xi$. Also, if $X$ is integral we have
$\mathcal{K}_{X, x} = R(X)$ is equal to the function field of $X$,
so $s/s_\xi \in R(X)$. If $s$ is a regular meromorphic section,
then actually $s/s_\xi \in R(X)^*$. On an integral scheme a regular
meromorphic section is the same thing as a nonzero meromorphic section.
Finally, we see that $s/s_\xi$ is independent of the choice of $s_\xi$ up to
multiplication by a unit of the local ring $\mathcal{O}_{X, x}$.
Putting everything together we see the following definition makes sense.

\begin{definition}
\label{definition-order-vanishing-meromorphic}
Let $X$ be a locally Noetherian integral scheme.
Let $\mathcal{L}$ be an invertible $\mathcal{O}_X$-module.
Let $s \in \Gamma(X, \mathcal{K}_X(\mathcal{L}))$
be a regular meromorphic section of $\mathcal{L}$.
For every prime divisor $Z \subset X$ we define the
{\it order of vanishing of $s$ along $Z$} as the integer
$$
\text{ord}_{Z, \mathcal{L}}(s)
= \text{ord}_{\mathcal{O}_{X, \xi}}(s/s_\xi)
$$
where the right hand side is the notion of
Algebra, Definition \ref{algebra-definition-ord},
$\xi \in Z$ is the generic point,
and $s_\xi \in \mathcal{L}_\xi$ is a generator.
\end{definition}

\noindent
As in the case of principal divisors we have the following lemma.

\begin{lemma}
\label{lemma-divisor-meromorphic-locally-finite}
Let $X$ be a locally Noetherian integral scheme. Let $\mathcal{L}$ be an
invertible $\mathcal{O}_X$-module. Let $s \in \mathcal{K}_X(\mathcal{L})$ be a
regular (i.e., nonzero) meromorphic section of $\mathcal{L}$. Then the sets
$$
\{Z \subset X \mid Z \text{ a prime divisor with generic point }\xi
\text{ and }s\text{ not in }\mathcal{L}_\xi\}
$$
and
$$
\{Z \subset X \mid Z \text{ is a prime divisor and }
\text{ord}_{Z, \mathcal{L}}(s) \not = 0\}
$$
are locally finite in $X$.
\end{lemma}

\begin{proof}
There exists a nonempty open subscheme $U \subset X$ such that $s$
corresponds to a section of $\Gamma(U, \mathcal{L})$ which generates
$\mathcal{L}$ over $U$. Hence the prime divisors which can occur
in the sets of the lemma are all irreducible components of $X \setminus U$.
Hence Lemma \ref{lemma-components-locally-finite}.
gives the desired result.
\end{proof}

\begin{lemma}
\label{lemma-divisor-meromorphic-well-defined}
Let $X$ be a locally Noetherian integral scheme.
Let $\mathcal{L}$ be an invertible $\mathcal{O}_X$-module.
Let $s, s' \in \mathcal{K}_X(\mathcal{L})$ be nonzero
meromorphic sections of $\mathcal{L}$. Then $f = s/s'$
is an element of $R(X)^*$ and we have
$$
\sum \text{ord}_{Z, \mathcal{L}}(s)[Z]
=
\sum \text{ord}_{Z, \mathcal{L}}(s')[Z]
+
\text{div}(f)
$$
as Weil divisors.
\end{lemma}

\begin{proof}
This is clear from the definitions.
Note that Lemma \ref{lemma-divisor-meromorphic-locally-finite}
guarantees that the sums are indeed Weil divisors.
\end{proof}

\begin{definition}
\label{definition-divisor-invertible-sheaf}
Let $X$ be a locally Noetherian integral scheme.
Let $\mathcal{L}$ be an invertible $\mathcal{O}_X$-module.
\begin{enumerate}
\item For any nonzero meromorphic section $s$ of $\mathcal{L}$
we define the {\it Weil divisor associated to $s$} as
$$
\text{div}_\mathcal{L}(s) =
\sum \text{ord}_{Z, \mathcal{L}}(s) [Z] \in \text{Div}(X)
$$
where the sum is over prime divisors.
\item We define {\it Weil divisor class associated to $\mathcal{L}$}
as the image of $\text{div}_\mathcal{L}(s)$ in $\text{Cl}(X)$
where $s$ is any nonzero meromorphic section of $\mathcal{L}$ over
$X$. This is well defined by
Lemma \ref{lemma-divisor-meromorphic-well-defined}.
\end{enumerate}
\end{definition}

\noindent
As expected this construction is additive in the invertible module.

\begin{lemma}
\label{lemma-c1-additive}
Let $X$ be a locally Noetherian integral scheme.
Let $\mathcal{L}$, $\mathcal{N}$ be invertible $\mathcal{O}_X$-modules.
Let $s$, resp.\ $t$ be a nonzero meromorphic section
of $\mathcal{L}$, resp.\ $\mathcal{N}$. Then $st$ is a nonzero
meromorphic section of $\mathcal{L} \otimes \mathcal{N}$, and
$$
\text{div}_{\mathcal{L} \otimes \mathcal{N}}(st)
=
\text{div}_\mathcal{L}(s) + \text{div}_\mathcal{N}(t)
$$
in $\text{Div}(X)$. In particular, the Weil divisor class of
$\mathcal{L} \otimes_{\mathcal{O}_X} \mathcal{N}$ is the sum
of the Weil divisor classes of $\mathcal{L}$ and $\mathcal{N}$.
\end{lemma}

\begin{proof}
Let $s$, resp.\ $t$ be a nonzero meromorphic section
of $\mathcal{L}$, resp.\ $\mathcal{N}$. Then $st$ is a nonzero
meromorphic section of $\mathcal{L} \otimes \mathcal{N}$.
Let $Z \subset X$ be a prime divisor. Let $\xi \in Z$ be its generic
point. Choose generators $s_\xi \in \mathcal{L}_\xi$, and
$t_\xi \in \mathcal{N}_\xi$. Then $s_\xi t_\xi$ is a generator
for $(\mathcal{L} \otimes \mathcal{N})_\xi$.
So $st/(s_\xi t_\xi) = (s/s_\xi)(t/t_\xi)$.
Hence we see that
$$
\text{ord}_{Z, \mathcal{L} \otimes \mathcal{N}}(st)
=
\text{ord}_{Z, \mathcal{L}}(s) + \text{ord}_{Z, \mathcal{N}}(t)
$$
by the additivity of the $\text{ord}_Z$ function.
\end{proof}

\noindent
In this way we obtain a homomorphism of abelian groups
\begin{equation}
\label{equation-c1}
\Pic(X) \longrightarrow \text{Cl}(X)
\end{equation}
which assigns to an invertible module its Weil divisor class.

\begin{lemma}
\label{lemma-normal-c1-injective}
Let $X$ be a locally Noetherian integral scheme. If $X$ is normal,
then the map (\ref{equation-c1}) $\Pic(X) \to \text{Cl}(X)$
is injective.
\end{lemma}

\begin{proof}
Let $\mathcal{L}$ be an invertible $\mathcal{O}_X$-module whose
associated Weil divisor class is trivial. Let $s$ be a regular
meromorphic section of $\mathcal{L}$. The assumption means that
$\text{div}_\mathcal{L}(s) = \text{div}(f)$ for some
$f \in R(X)^*$. Then we see that $t = f^{-1}s$ is a regular
meromorphic section of $\mathcal{L}$ with
$\text{div}_\mathcal{L}(t) = 0$, see
Lemma \ref{lemma-divisor-meromorphic-well-defined}.
We will show that $t$ defines a trivialization of $\mathcal{L}$
which finishes the proof of the lemma.
In order to prove this we may work locally on $X$.
Hence we may assume that $X = \Spec(A)$ is affine
and that $\mathcal{L}$ is trivial. Then $A$ is a Noetherian normal
domain and $t$ is an element of its fraction field
such that $\text{ord}_{A_\mathfrak p}(t) = 0$
for all height $1$ primes $\mathfrak p$ of $A$.
Our goal is to show that $t$ is a unit of $A$.
Since $A_\mathfrak p$ is a discrete valuation ring for height
one primes of $A$ (Algebra, Lemma \ref{algebra-lemma-criterion-normal}), the
condition signifies that $t \in A_\mathfrak p^*$ for all primes $\mathfrak p$
of height $1$. This implies $t \in A$ and $t^{-1} \in A$ by
Algebra, Lemma
\ref{algebra-lemma-normal-domain-intersection-localizations-height-1}
and the proof is complete.
\end{proof}

\begin{lemma}
\label{lemma-local-rings-UFD-c1-bijective}
Let $X$ be a locally Noetherian integral scheme. Consider the map
(\ref{equation-c1}) $\Pic(X) \to \text{Cl}(X)$.
The following are equivalent
\begin{enumerate}
\item the local rings of $X$ are UFDs, and
\item $X$ is normal and $\Pic(X) \to \text{Cl}(X)$
is surjective.
\end{enumerate}
In this case $\Pic(X) \to \text{Cl}(X)$ is an isomorphism.
\end{lemma}

\begin{proof}
If (1) holds, then $X$ is normal by
Algebra, Lemma \ref{algebra-lemma-UFD-normal}.
Hence the map (\ref{equation-c1}) is injective by
Lemma \ref{lemma-normal-c1-injective}. Moreover,
every prime divisor $D \subset X$ is an effective
Cartier divisor by Lemma \ref{lemma-weil-divisor-is-cartier-UFD}.
In this case the canonical section $1_D$ of $\mathcal{O}_X(D)$
(Definition \ref{definition-invertible-sheaf-effective-Cartier-divisor})
vanishes exactly along $D$ and we see that the class of $D$ is the
image of $\mathcal{O}_X(D)$ under the map (\ref{equation-c1}).
Thus the map is surjective as well.

\medskip\noindent
Assume (2) holds. Pick a prime divisor $D \subset X$.
Since (\ref{equation-c1}) is surjective there exists an invertible
sheaf $\mathcal{L}$, a regular meromorphic section $s$, and $f \in R(X)^*$
such that $\text{div}_\mathcal{L}(s) + \text{div}(f) = [D]$.
In other words, $\text{div}_\mathcal{L}(fs) = [D]$.
Let $x \in X$ and let $A = \mathcal{O}_{X, x}$. Thus $A$ is
a Noetherian local normal domain with fraction field $K = R(X)$.
Every height $1$ prime of $A$ corresponds to a prime divisor on $X$
and every invertible $\mathcal{O}_X$-module restricts to the
trivial invertible module on $\Spec(A)$. It follows that for every
height $1$ prime $\mathfrak p \subset A$ there exists an element $f \in K$
such that $\text{ord}_{A_\mathfrak p}(f) = 1$ and
$\text{ord}_{A_{\mathfrak p'}}(f) = 0$ for every other
height one prime $\mathfrak p'$. Then $f \in A$ by Algebra, Lemma
\ref{algebra-lemma-normal-domain-intersection-localizations-height-1}.
Arguing in the same fashion we see that every element $g \in \mathfrak p$
is of the form $g = af$ for some $a \in A$. Thus we see that every
height one prime ideal of $A$ is principal and $A$ is a UFD
by Algebra, Lemma \ref{algebra-lemma-characterize-UFD-height-1}.
\end{proof}







